\documentclass{IEEEtran}
\usepackage{graphicx}
\usepackage{booktabs}
\usepackage{amsmath}
\usepackage{amssymb}
\usepackage[hidelinks]{hyperref}
\usepackage{url}
\usepackage[numbers,sort&compress]{natbib}
\title{MP-CycleGAN: Multi-Perceptual Constraints for Unpaired Underwater Image Enhancement}
\author{Anonymous Authors}
\begin{document}
\maketitle
\begin{abstract}
Underwater image enhancement (UIE) is a critical preprocessing step for marine robotics and underwater vision. Real underwater images often suffer from wavelength-dependent absorption and scattering, leading to color casts, low contrast, and haze-like artifacts. While supervised UIE models can achieve strong full-reference metrics, they rely on pixel-aligned training pairs that are difficult to acquire and may generalize poorly across domains. This paper presents MP-CycleGAN, an unpaired UIE framework built on CycleGAN with multi-perceptual constraints. Here, multi-perceptual (MP) denotes three complementary perceptual perspectives that jointly regularize the unpaired mapping: color perception via a gray-world prior, structure perception via gradient consistency, and semantic perception via a VGG-based perceptual loss. We evaluate the method under a unified inference and evaluation pipeline on EUVP (in-domain) and UIEB (cross-domain), and report bootstrap confidence intervals for PSNR/SSIM to improve reproducibility. Results show small but consistent gains in matched in-domain evaluation and competitive cross-domain performance with improved perceptual-quality tendencies.
\end{abstract}
\begin{IEEEkeywords}
Underwater image enhancement, unpaired learning, CycleGAN, perceptual loss, gray-world prior, structure consistency.
\end{IEEEkeywords}
\section{Introduction}
Underwater images are widely used in marine exploration, inspection, ecological monitoring, and robot navigation. Due to wavelength-dependent absorption and scattering, underwater images frequently exhibit color casts and reduced visibility, which degrade both human interpretation and downstream vision tasks. UIE aims to improve perceptual quality while preserving structural content.
Collecting pixel-aligned underwater/clean pairs is challenging in real environments. Unpaired image-to-image translation frameworks such as CycleGAN alleviate paired data requirements, but may introduce color drift, structural distortions, and hallucinated textures. This paper proposes MP-CycleGAN to stabilize unpaired UIE by constraining the mapping using complementary perceptual cues.
\section{Method}
\subsection{CycleGAN Base}
We consider underwater domain $A$ and enhanced domain $B$. CycleGAN learns generators $G: A\rightarrow B$ and $F: B\rightarrow A$ with discriminators $D_B$ and $D_A$, using adversarial and cycle-consistency objectives:
\begin{equation}
\mathcal{L}_{total}=\mathcal{L}_{GAN}+\lambda_{cyc}\mathcal{L}_{cyc}+\lambda_{idt}\mathcal{L}_{idt}+\mathcal{L}_{MP}.
\end{equation}
\subsection{Multi-Perceptual Constraints}
\begin{equation}
\mathcal{L}_{MP}=\lambda_{gray}\mathcal{L}_{gray}+\lambda_{struct}\mathcal{L}_{struct}+\lambda_{perc}\mathcal{L}_{perc}+\lambda_{color}\mathcal{L}_{color}.
\end{equation}
\subsubsection{Gray-world color prior}
\begin{equation}
\mathcal{L}_{gray}(\hat{I})=\frac{1}{3}\sum_{c\in\{r,g,b\}}\left|\mu_c(\hat{I})-\mu_{gray}(\hat{I})\right|,\quad
\mu_{gray}(\hat{I})=\frac{\mu_r(\hat{I})+\mu_g(\hat{I})+\mu_b(\hat{I})}{3}.
\end{equation}
\subsubsection{Gradient-based structure consistency}
\begin{equation}
\mathcal{L}_{struct}(I,\hat{I})=\left\|\nabla_x \hat{I}-\nabla_x I\right\|_1+\left\|\nabla_y \hat{I}-\nabla_y I\right\|_1.
\end{equation}
\subsubsection{VGG perceptual loss}
\begin{equation}
\mathcal{L}_{perc}(I,\hat{I})=\left\|\phi(I)-\phi(\hat{I})\right\|_1.
\end{equation}
\subsubsection{Optional global color statistics}
\begin{equation}
\mathcal{L}_{color}(I,\hat{I})=\left\|\mu(\hat{I})-\mu(I)\right\|_1+\left\|\sigma(\hat{I})-\sigma(I)\right\|_1.
\end{equation}
\section{Experiments}
\subsection{Datasets and Protocol}
We evaluate on EUVP (matched=200) for in-domain testing and UIEB (matched=890) for cross-domain testing without fine-tuning on UIEB.
\subsection{Metrics}
We report PSNR/SSIM and no-reference metrics UCIQE/UIQM. We compute 95\% bootstrap confidence intervals for PSNR/SSIM with \texttt{bootstrap\_iters=2000} and \texttt{seed=123}.
\subsection{Reproducibility}
For reproducibility, the final submission should provide the exact training configuration (loss weights, perceptual layer), inference outputs for matched subsets, and evaluation scripts used to compute all metrics.
\subsection{Downstream Object Detection}
We additionally evaluate downstream object detection on the URPC optical object detection dataset~\cite{urpc2020} using a YOLO-based detector. We compare detection performance on raw validation images versus MP-CycleGAN enhanced validation images under identical labels, reporting mAP@0.5 and mAP@0.5:0.95.
When using a detector trained only on raw images, MP-CycleGAN enhancement reduces detection accuracy, suggesting distribution shift. Fine-tuning the detector on MP-CycleGAN enhanced training images recovers performance on enhanced inputs, while mixed training on raw and MP-CycleGAN enhanced data improves robustness across both domains.
\begin{figure}[!t]
\centering
\includegraphics[width=\linewidth]{../docs/figures/downstream_detection_fixed_vs_mixed.png}
\caption{Qualitative downstream detection comparison on URPC (top: fixed detector trained on raw; bottom: mixed detector trained on raw and MP-CycleGAN enhanced images). Green: ground-truth boxes. Red: predicted boxes.}
\end{figure}
\begin{table}[!t]
\centering
\caption{Downstream object detection performance on URPC (validation split).}
\begin{tabular}{lcc}
\toprule
Setting & mAP@0.5 & mAP@0.5:0.95 \\
\midrule
Fixed detector (raw val) & 0.8365 & 0.4908 \\
Fixed detector (MP-CycleGAN val) & 0.4614 & 0.2157 \\
Enhancement-aware detector (raw val) & 0.5787 & 0.2672 \\
Enhancement-aware detector (MP-CycleGAN val) & 0.6261 & 0.3181 \\
Mixed detector (raw val) & $0.8458\pm0.0005$ & $0.4982\pm0.0011$ \\
Mixed detector (MP-CycleGAN val) & $0.6601\pm0.0034$ & $0.3388\pm0.0003$ \\
\bottomrule
\end{tabular}
\end{table}
\section{Results}
\subsection{In-domain Results on EUVP}
\begin{table}[!t]
\centering
\caption{In-domain results on EUVP matched test set (matched=200).}
\begin{tabular}{lcccc}
\toprule
Method & PSNR (dB) & SSIM & UCIQE & UIQM \\
\midrule
Identity & 20.33 & 0.794 & 24.88 & 5.77 \\
Gray-World & 19.42 & 0.784 & 24.27 & 6.72 \\
CycleGAN & 22.82 & 0.783 & 25.93 & 5.33 \\
MP-CycleGAN & \textbf{22.91} & \textbf{0.795} & 25.82 & 5.41 \\
\bottomrule
\end{tabular}
\end{table}
\subsection{Cross-domain Results on UIEB}
\begin{table}[!t]
\centering
\caption{Cross-domain testing on UIEB matched subset (matched=890).}
\begin{tabular}{lccc}
\toprule
Metric & Input & CycleGAN & MP-CycleGAN \\
\midrule
PSNR (dB) & - & \textbf{17.84} & 17.54 \\
SSIM & - & \textbf{0.626} & 0.624 \\
UCIQE & 21.70 & \textbf{23.25} & 22.46 \\
UIQM & 6.80 & 10.48 & \textbf{10.50} \\
\bottomrule
\end{tabular}
\end{table}
\subsection{Qualitative Results}
\begin{figure}[!t]
\centering
\includegraphics[width=\linewidth]{../docs/figures/visual_comparison_final_refined.png}
\caption{Qualitative comparison on EUVP (input, CycleGAN, MP-CycleGAN, reference).}
\end{figure}
\section{Conclusion}
We presented MP-CycleGAN for unpaired underwater image enhancement. By integrating multi-perceptual constraints spanning color, structure, and semantic consistency, the method improves stability of unpaired translation for underwater enhancement. Experiments on EUVP and UIEB show small but consistent in-domain gains and competitive cross-domain performance with favorable perceptual-quality tendencies.
\begin{thebibliography}{13}
\bibitem{li2020uieb} C. Li, C. Guo, W. Ren, R. Cong, J. Hou, S. Kwong, and D. Tao, ``An Underwater Image Enhancement Benchmark Dataset and Beyond,'' \emph{IEEE Transactions on Image Processing}, vol. 29, pp. 4376--4389, 2020.
\bibitem{islam2020funiegan} M. J. Islam, Y. Xia, and J. Sattar, ``Fast Underwater Image Enhancement for Improved Visual Perception,'' \emph{IEEE Robotics and Automation Letters}, vol. 5, no. 2, pp. 3227--3234, 2020.
\bibitem{li2021ucolor} C. Li, S. Anwar, J. Hou, R. Cong, C. Guo, and W. Ren, ``Underwater Image Enhancement via Medium Transmission-Guided Multi-Color Space Embedding,'' \emph{IEEE Transactions on Image Processing}, vol. 30, pp. 4985--5000, 2021.
\bibitem{zhu2017cyclegan} J.-Y. Zhu, T. Park, P. Isola, and A. A. Efros, ``Unpaired Image-to-Image Translation Using Cycle-Consistent Adversarial Networks,'' in \emph{Proc. ICCV}, 2017, pp. 2242--2251.
\bibitem{wang2004ssim} Z. Wang, A. C. Bovik, H. R. Sheikh, and E. P. Simoncelli, ``Image Quality Assessment: From Error Visibility to Structural Similarity,'' \emph{IEEE Transactions on Image Processing}, vol. 13, no. 4, pp. 600--612, 2004.
\bibitem{yang2015uciqe} M. Yang and A. Sowmya, ``An Underwater Color Image Quality Evaluation Metric,'' \emph{IEEE Transactions on Image Processing}, vol. 24, no. 12, pp. 6062--6071, 2015.
\bibitem{panetta2016uiqm} K. Panetta, C. Gao, and S. Agaian, ``Human-Visual-System-Inspired Underwater Image Quality Measures,'' \emph{IEEE Journal of Oceanic Engineering}, vol. 41, no. 3, pp. 541--551, 2016.
\bibitem{johnson2016perceptual} J. Johnson, A. Alahi, and L. Fei-Fei, ``Perceptual Losses for Real-Time Style Transfer and Super-Resolution,'' in \emph{Proc. ECCV}, 2016.
\bibitem{simonyan2015vgg} K. Simonyan and A. Zisserman, ``Very Deep Convolutional Networks for Large-Scale Image Recognition,'' in \emph{Proc. ICLR}, 2015.
\bibitem{buchsbaum1980grayworld} G. Buchsbaum, ``A Spatial Processor Model for Object Colour Perception,'' \emph{Journal of the Franklin Institute}, vol. 310, no. 1, pp. 1--26, 1980.
\bibitem{urpc2020} URPC, ``Underwater Robot Picking Competition (Optical Image Object Detection) Dataset,'' 2020. [Online]. Available: \url{https://openi.pcl.ac.cn/OpenOrcinus_orca/URPC_opticalimage_dataset/datasets}.
\end{thebibliography}
\end{document}
